\subsection{正线性泛函}

定义 2. --- 设 A 是对合代数。A 上的线性泛函 \(\lambda\) 称为正的，如果 \(\lambda(a^{*}a) \in \mathbf{R}_{+}\) 对所有 \(a \in \mathrm{A}\) 都成立。

若 A 是对合巴拿赫代数，我们以 \(A_{+}^{\prime}\) 表示 A 上连续正线性泛函的集合。

设 A 为对合巴拿赫代数。集合 \(A_{+}^{\prime}\) 是 A 上 C-线性泛函所在实向量空间中的尖凸锥。

引理 2. --- 设 A 为对合巴拿赫代数，且 \(\lambda\) 是 A 上的正线性泛函。

\begin{enumerate}
\def\labelenumi{\alph{enumi})}
\tightlist
\item
  对 A 中任意 \(a\) 和 \(b\)，有 \(\lambda(a^{*}b)=\overline{\lambda(b^{*}a)}\)，且
\end{enumerate}

\[
\left| \lambda \left(a ^ {*} b\right) \right| ^ {2} \leqslant \lambda \left(a ^ {*} a\right) \lambda \left(b ^ {*} b\right);
\]

\begin{enumerate}
\def\labelenumi{\alph{enumi})}
\setcounter{enumi}{1}
\tightlist
\item
  若 A 含单位元，则线性泛函 \(\lambda\) 连续，且其范数等于 \(\lambda(1)\)。
\end{enumerate}

映射 \((a,b)\mapsto\lambda(a^{*}b)\) 是 A 上的正埃尔米特形式；因此，对 A 中任意 a 和 b，它满足 \(\lambda(a^{*}b)=\overline{\lambda(b^{*}a)}\) 以及 \(|\lambda(a^{*}b)|^{2}\leqslant\lambda(a^{*}a)\lambda(b^{*}b)\)（EVT, V, p. 2, 注以及 p. 3, 命题 2）。

证明 \(b\)。设 \(a \in A\) 是范数 \(< 1\) 的埃尔米特元。\(a\) 的谱半径小于 \(\|a\|\)，所以 \(a\) 的谱包含于 C 的开单位圆盘 D 中（I, p. 24, 定理 1）。设 \(f\) 是 D 上满足 \(f(z)^{2} = 1 - z\) 且对所有 \(z \in D\) 成立的全纯函数（V, p. 435, 引理 1）。对元 \(a\) 和函数 \(f\) 应用全纯函数演算。元 \(b = f(a)\) 满足 \(b^{2} = 1 - a\)（I, p. 74, 定理 5）。由 V, p. 435, 命题 3，还有 \(f(a)^{*} = f^{*}(a^{*}) = f(a)\)，故 \(b\) 是埃尔米特元。于是 \(\lambda(1 - a) = \lambda(b^{*}b) \geqslant 0\)，从而 \(\lambda(a) \leqslant \lambda(1)\)。

现在设 \(b \in A\) 且范数 \(< 1\)。元 \(b^*b\) 是埃尔米特元且 \(\|b^*b\| < 1\)，因此在应用 \(a\) )（取 \(a = 1\)）的情形，得到

\[
| \lambda (b) | ^ {2} \leqslant \lambda (1) \lambda (b ^ {*} b) \leqslant \lambda (1) ^ {2},
\]

当 \(b = 1\) 时取等号。因此线性泛函 \(\lambda\) 连续，且其范数等于 \(\lambda(1)\)。断言 \(b\) ) 得证。

例。--- 设 X 是紧拓扑空间，且 A 是星代数 \(\mathcal{C}(\mathrm{X})\)。A 上的正线性泛函与 X 上的正测度相对应（INT, III, p. 52, § 1, n° 6, 定理 1）。

引理 3. --- 设 A 是具有逼近单位元的含单位元对合巴拿赫代数（I, p. 120, 定义 7）。设 λ 是 A 上的连续正线性泛函。

\begin{enumerate}
\def\labelenumi{\alph{enumi})}
\item
  对 A 中任意 a，有 \(\lambda(a^{*})=\overline{\lambda(a)}\) 以及 \(|\lambda(a)|^{2}\leqslant\|\lambda\|\lambda(a^{*}a)\)；
\item
  设 \(\widetilde{\mathrm{A}}\) 是通过向 A 添加单位元得到的对合代数，e 是其单位元。存在 \(\widetilde{\lambda}\) 在 \(\widetilde{\mathrm{A}}\) 上的连续正线性泛函，它延拓 \(\lambda\)，且满足 \(\widetilde{\lambda}(e) = \|\lambda\|\)；
\item
  对 A 中任意 \(a\) 和 \(b\)，有 \(|\lambda(b^{*}ab)|\leqslant\|a\|\lambda(b^{*}b)\)。
\end{enumerate}

证明 \(a\) )。设 \(\mathfrak{F}\) 是 A 的逼近单位元，且 \(a \in \mathrm{A}\)。利用引理 2、\(a\) ) 以及逼近单位元的定义，可得

\[
\lambda (a ^ {*}) = \lim _ {f, \mathfrak {F}} \lambda (f a ^ {*}) = \lim _ {f, \mathfrak {F}} \overline {{\lambda (a f ^ {*})}} = \lim _ {f, \mathfrak {F}} \overline {{\lambda ((f a ^ {*}) ^ {*})}} = \overline {{\lambda (a)}},
\]

因此（同上）

\[
| \lambda (a) | ^ {2} = \lim _ {f, \mathfrak {F}} | \lambda (f a) | ^ {2} \leqslant \lambda (a ^ {*} a) \operatorname * {limsup} _ {f, \mathfrak {F}} \lambda (f f ^ {*}) \leqslant \| \lambda \| \lambda (a ^ {*} a),
\]

因为 \(\mathfrak{F}\) 是 A 的单位球上的滤子。

证明 \(b\) )。不妨设 \(\lambda\) 非零，进而设 \(\| \lambda \| = 1\)。对 \(a \in \mathrm{A}\) 和 \(z \in \mathbf{C}\)，定义 \(\widetilde{\lambda}(a + z \cdot e) = \lambda(a) + z\)。映射 \(\widetilde{\lambda}\) 是 \(\widetilde{\mathrm{A}}\) 上延拓 \(\lambda\) 的连续线性泛函，并满足 \(\widetilde{\lambda}(e) = 1\)。它是正的：对任意 \(a \in \mathrm{A}\) 和 \(z \in \mathbf{C}\)，根据 \(a\) )，计算得

\[
\begin{array}{r l} \widetilde {\lambda} ((a + z \cdot e) ^ {*} (a + z \cdot e)) & = \lambda (a ^ {*} a) + \overline {{z}} \lambda (a) + z \lambda (a ^ {*}) + | z | ^ {2} \\ & = | z + \lambda (a) | ^ {2} + \lambda (a ^ {*} a) - | \lambda (a) | ^ {2} \geqslant 0. \end{array}
\]

最后，设 \(b\) 是 A 的一个元。映射 \(a \mapsto \widetilde{\lambda}(b^{*}ab)\) 是含单位元对合巴拿赫代数 \(\widetilde{A}\) 上的正线性泛函。因此它连续，且范数等于其在 \(e\) 处的值（引理 2，\(b\) )），而该值等于 \(\lambda(b^{*}b)\)，从而得到断言 \(c\) )。

命题 4. — 设 A 为对合巴拿赫代数。

\begin{enumerate}
\def\labelenumi{\alph{enumi})}
\item
  对每个希尔伯特空间 E、每个对合代数态射 \(\varphi\colon A\to\mathcal{L}(E)\) 以及每个向量 \(x\in E\)，在 A 上由 \(\lambda(a)=\langle x|\varphi(a)x\rangle\) 定义的线性泛函是连续正线性泛函；
\item
  设 \(\lambda \in \mathrm{A}_{+}^{\prime}\)。映射 \(f\) 从 \(\mathrm{A} \times \mathrm{A}\) 到 \(\mathbf{C}\)，由 \(f(a, b) = \lambda (a^{*}b)\) 对所有 \((a, b) \in \mathrm{A} \times \mathrm{A}\) 定义，是 \(\mathrm{A}\) 上的普遍正核。
\end{enumerate}

证明 \(a\) )。对每个 \(a \in \mathrm{A}\)，有

\[
\lambda (a ^ {*} a) = \langle x \mid \varphi (a ^ {*} a) x \rangle = \| \varphi (a) x \| ^ {2} \geqslant 0
\]

所以 \(\lambda\) 是 A 上的正线性泛函；由于 \(\varphi\) 连续，它也是连续的（I, p. 104, 命题 2）。

证明 \(b\) )。函数 \(f\) 连续。对每个 \(n \in \mathbf{N}\)、A 中的每个族 \((a_i)_{0 \leqslant i \leqslant n}\) 以及每个复数族 \((t_i)_{0 \leqslant i \leqslant n}\)，有

\[
\sum_ {i = 0} ^ {n} \sum_ {j = 0} ^ {n} \bar {t} _ {i} t _ {j} f (a _ {i}, a _ {j}) = \lambda \left(\left(\sum_ {i = 0} ^ {n} t _ {i} a _ {i}\right) ^ {*} \left(\sum_ {j = 0} ^ {n} t _ {j} a _ {j}\right)\right) \geqslant 0,
\]

从而得到结论（V, p. 433, 定义 1）。

定义 3. — 设 A 是对合巴拿赫代数，且 \(\lambda \in \mathrm{A}_{+}^{\prime}\) 是 A 上的连续正线性泛函。\(\lambda\) 的希尔伯特实现是一个三元组 (E, \(x, \varphi\) )，其中 E 是希尔伯特空间，\(x\) 是 E 的一个元，\(\varphi\) 是从 A 到 \(\mathcal{L}(\mathrm{E})\) 的对合代数态射，并且 \(\lambda(b^{*}a) = \langle \varphi(b)x | \varphi(a)x \rangle\) 对所有 \((a,b) \in \mathrm{A}^{2}\) 都成立。

如果元素 \(\varphi(a)x\)（其中 \(a\in\mathrm{A}\)）在 E 中构成全集，则称 (E, \(x,\varphi\) ) 是 \(\lambda\) 的循环希尔伯特空间实现。

命题 5. --- 设 A 是具有逼近单位元的对合巴拿赫代数。A 上的每个连续正线性泛函都有一个循环希尔伯特实现 (E, x, φ)。若 A 含单位元，则存在一个其中态射 φ 保单位元的此类实现。

由引理 3，b)，不妨设 A 是含单位元的代数。

设 \(\lambda\) 是 A 上的连续正线性泛函。设 (E, \(g\) ) 是 A 上的普遍正核 \(f\) 的循环希尔伯特实现，该核由 \(f(a,b) = \lambda(a^{*}b)\) 定义（命题 4 和定理 1）。映射 \(g\) 连续，其像在 E 中是全的，并且 \(\lambda(a^{*}b) = \langle g(a)|g(b)\rangle\) 对所有 \((a,b) \in \mathrm{A}^2\) 都成立。置 \(x = g(1) \in \mathrm{E}\)。

从 A 到 E 的映射 g 是线性的：事实上，对 \((a,b,c)\in\mathrm{A}^{3}\) 和 \((s,t)\in\mathbf{C}^{2}\)，有

\[
\begin{array}{c} \langle g (c) | g (s a + t b) \rangle = \lambda (c ^ {*} (s a + t b)) \\ = s \lambda (c ^ {*} a) + t \lambda (c ^ {*} b) = \langle g (c) | s g (a) + t g (b) \rangle , \end{array}
\]

由于 \(g\) 的像在 E 中是全的，故断言成立。

特别地，g 的像 F 是 E 的稠密向量子空间。g 的核是 A 的左理想：若 \(g(b) = 0\)，则对 A 中任意 a 和 c，有

\[
\langle g (a b) \mid g (c) \rangle = \lambda ((a b) ^ {*} c) = \lambda (b ^ {*} (a ^ {*} c)) = \langle g (b) \mid g (a ^ {*} c) \rangle = 0,
\]

由于 F 在 E 中稠密，所以 \(g(ab) = 0\)。

设 \(a \in \mathrm{A}\)。由于 \(g\) 的核是 \(\mathrm{A}\) 的左理想，存在一个线性映射 \(\widetilde{\varphi}(a) \colon \mathrm{F} \to \mathrm{F}\)，使得 \(\widetilde{\varphi}(a)(g(b)) = g(ab)\) 对所有 \(b \in \mathrm{A}\) 成立。此外，对每个 \(b \in \mathrm{A}\)，有

\[
\| \widetilde {\varphi} (a) (g (b)) \| ^ {2} = \| g (a b) \| ^ {2} = \lambda \left(b ^ {*} a ^ {*} a b\right) \leqslant \| a ^ {*} a \| \lambda \left(b ^ {*} b\right) \leqslant \| a \| ^ {2} \| g (b) \| ^ {2}
\]

（引理 3，c)），所以 \(\widetilde{\varphi}(a)\) 连续且范数 \(\leqslant \|a\|\)。因此存在唯一的自同态 \(\varphi(a) \in \mathcal{L}(\mathrm{E})\)，它通过对子空间的传递诱导出 \(\widetilde{\varphi}(a)\)。

设 \(a\) 和 \(b\) 属于 A。有

\[
(\varphi (a) \circ \varphi (b)) (g (c)) = g (a b c) = \varphi (a b) (g (c))
\]

对 A 中任意 c 都成立，因此由于 F 在 E 中稠密，\(\varphi(ab) = \varphi(a) \circ \varphi(b)\)。

映射 \(\varphi\) 从 A 到 \(\mathcal{L}(\mathrm{E})\) 是线性的：事实上，对任意 \((a,b)\in\mathrm{A}^{2}\) 和 \((s,t)\in\mathbf{C}^{2}\)，有

\[
\varphi (s a + t b) (g (c)) = g ((s a + t b) c) = (s \varphi (a) + t \varphi (b)) g (c)
\]

对 A 中任意 \(c \in A\) 都成立，从而由于 F 在 E 中稠密，\(\varphi(sa + tb) = s\varphi(a) + t\varphi(b)\)。由于 \(\|\varphi(b)\| \leqslant \|b\|\) 对所有 \(b \in A\) 成立，映射 \(\varphi\) 连续。又有 \(\varphi(1) = 1_{\mathrm{E}}\)，因为 \(\varphi(1)(g(a)) = g(a)\) 对所有 \(a \in A\) 成立，且 F 在 E 中稠密。

最后，设 a、b、c 属于 A。有

\[
\begin{array}{c} \langle g (c) \mid \varphi (a ^ {*}) (g (b)) \rangle = \langle g (c) \mid g (a ^ {*} b) \rangle = \lambda (c ^ {*} a ^ {*} b) \\ = \langle g (a c) \mid g (b) \rangle = \langle \varphi (a) (g (c)) \mid g (b) \rangle \end{array}
\]

由于 F 在 E 中稠密，从而 \(\varphi(a^{*})=\varphi(a)^{*}\)。

总之，映射 \(\varphi\) 是从 A 到 \(\mathcal{L}(\mathrm{E})\) 的连续对合代数态射，且 \(g(a) = \varphi(a)x\) 对所有 \(a \in \mathrm{A}\) 成立；因此三元组 \((\mathrm{E}, x, \varphi)\) 是 \(\lambda\) 的循环希尔伯特实现。

命题 6. --- 设 A 是对合巴拿赫代数，λ 是 A 上的连续正线性泛函，且 (E₁, x₁, φ₁) 和 (E₂, x₂, φ₂) 是 λ 的希尔伯特实现。设 (E₁, x₁, φ₁) 是循环希尔伯特实现。

\begin{enumerate}
\def\labelenumi{\alph{enumi})}
\item
  存在唯一的从 \(E_{1}\) 到 \(E_{2}\) 的连续线性映射 u，它是从 \(\varphi_{1}\) 到 \(\varphi_{2}\) 的表示态射（I, p. 11），并且满足 \(u(\varphi_{1}(a)x_{1}) = \varphi_{2}(a)x_{2}\) 对所有 \(a \in A\) 成立；
\item
  线性映射 u 是等距的；
\item
  若 \(\left(\mathrm{E}_{2},x_{2},\varphi_{2}\right)\) 是循环的，则 u 是同构；
\item
  若 \(\left(\mathrm{E}_2,x_2,\varphi_2\right)\) 是循环的且 A 含单位元，则 \(u(x_{1}) = x_{2}\)。
\end{enumerate}

对 \(j \in \{1,2\}\)，定义 \(\gamma_j \colon A \to E_j\)，使得 \(\gamma_j(a) = \varphi_j(a)x_j\) 对所有 \(a \in A\) 成立。按定义，\((E_j, \gamma_j)\) 是由普遍正核 \(f\) 在 \(A\) 上定义的希尔伯特实现，该核由 \((a,b) \mapsto \lambda(a^*b)\) 定义。希尔伯特实现 \((E_1, \gamma_1)\) 是循环的；由 V, p. 434, 命题 1，存在唯一的连续线性映射 \(u \colon E_1 \to E_2\)，使得 \(\gamma_2 = u \circ \gamma_1\)，且该映射是等距的。此外，若 \((E_2, x_2, \varphi_2)\) 也是循环的，则 \(u\) 是同构。要证明 \(a)\)、\(b)\) 和 \(c)\)，只需证明 \(u\) 是从 \(\varphi_1\) 到 \(\varphi_2\) 的表示态射。

设 \(a \in \mathrm{A}\)。对每个 \(b \in \mathrm{A}\)，有

\[
\begin{array}{c} (u \circ \varphi_ {1} (a)) (\gamma_ {1} (b)) = (u \circ \varphi_ {1} (a b)) x _ {1} = (u \circ \gamma_ {1}) (a b) \\ = \gamma_ {2} (a b) = \varphi_ {2} (a) (\gamma_ {2} (b)) = \varphi_ {2} (a) (u (\gamma_ {1} (b))), \end{array}
\]

因此，连续线性映射 \(u \circ \varphi_1(a)\) 和 \(\varphi_2(a) \circ u\) 在 \(E_1\) 中由 \(\gamma_1\) 的像生成的子空间上一致；由假设，该子空间在 \(E_1\) 中稠密，所以 \(u \circ \varphi_1(a) = \varphi_2(a) \circ u\)。

最后证明 \(d\) )。因此设 A 含单位元，且 \((\mathrm{E}_2, x_2, \varphi_2)\) 是循环的。存在循环希尔伯特空间实现 \((\mathrm{E}_3, x_3, \varphi_3)\) 的 \(\lambda\)，使得 \(\varphi_3\) 是保单位元的态射（命题 5）。设 \(j \in \{1, 2\}\)。将前述结果应用于 \((\mathrm{E}_j, x_j, \varphi_j)\) 和 \((\mathrm{E}_3, x_3, \varphi_3)\)，可知存在等距同构 \(\widetilde{u}_j\)，从 \(\mathrm{E}_j\) 到 \(\mathrm{E}_3\)，使得 \(\widetilde{u}_j \circ \varphi_j(a) = \varphi_3(a) \circ \widetilde{u}_j\) 对所有 \(a \in \mathrm{A}\) 成立。取 \(a = 1_{\mathrm{A}}\)，可知 \(\varphi_j\) 保单位元。于是 \(x_2 = \varphi_2(1_{\mathrm{A}})x_2 = u(\varphi_1(1_{\mathrm{A}})x_1) = u(x_1)\)。

注。--- 保持命题中的记号。可能 \(u(x_{1})\) 不同于 \(x_{2}\)（V, p. 493, 习题 3，b)）。不过，三元组 \((\mathrm{E}_{2}, u(x_{1}), \varphi_{2})\) 也是 \(\lambda\) 的希尔伯特实现。

